\section{Datensatz und Vorverarbeitung}
\label{sec:datensatz}

\subsection{Wahl des Datensatzes}

\begin{hinweisbox}[Platzhalter]
Hier beschreiben, welcher Datensatz verwendet wurde (z.\,B.\ XOR, Iris, MNIST, ein selbst erzeugter Datensatz) und warum die Wahl auf diesen Datensatz gefallen ist.
\end{hinweisbox}

\subsection{Struktur der Daten}

\begin{hinweisbox}[Platzhalter]
Angaben zu: Anzahl der Beispiele, Anzahl der Merkmale, Anzahl der Klassen, Format der Eingaben (z.\,B.\ Vektor der Länge $n$), Format der Labels (z.\,B.\ One-Hot).
\end{hinweisbox}

\subsection{Vorverarbeitung}

\begin{hinweisbox}[Platzhalter]
Schritte der Vorverarbeitung: Normalisierung (z.\,B.\ auf $[0,1]$ oder Mittelwert 0, Varianz 1), Flatten von Bildern, eventuell One-Hot-Encoding der Labels, Shuffling.
\end{hinweisbox}

\subsection{Trainings-/Validierungs-/Testaufteilung}

\begin{table}[htbp]
    \centering
    \small
    \begin{tabular}{l c c}
    \toprule
    \textbf{Teilmenge} & \textbf{Anteil} & \textbf{Anzahl Beispiele} \\
    \midrule
    Training    & 70\,\% & [HIER EINFÜGEN] \\
    Validierung & 15\,\% & [HIER EINFÜGEN] \\
    Test        & 15\,\% & [HIER EINFÜGEN] \\
    \bottomrule
    \end{tabular}
    \caption{Aufteilung des Datensatzes.}
    \label{tab:datenaufteilung}
\end{table}

\newpage